% =========================================================
% EDITE AQUI - DADOS DO TRABALHO COM IDENTIFICAÇÃO
% =========================================================
% Troque somente os textos dentro das chaves {...}.

% ATENÇÃO: o regulamento exige que o título seja EXATAMENTE igual ao do plano
% de trabalho cadastrado no SISPEQ pelo orientador. O título abaixo foi
% adaptado do artigo do SBrT e PRECISA ser conferido contra o SISPEQ antes
% da submissão -- divergência obriga reenvio.
\renewcommand{\sicitedadostitulo}{COMUNICAÇÃO SEMÂNTICA E APRENDIZADO FEDERADO NA BORDA: UMA ARQUITETURA BASEADA EM AUTOENCODERS PARA REDES 6G}

% Ordem exigida: apresentador(a), coautores(as) e, por último, o orientador.
\renewcommand{\sicitedadosautores}{Gian Pedro Rodrigues\textsuperscript{1*}, Herman Lima dos Santos\textsuperscript{1}}

\renewcommand{\sicitedadosafiliacoes}{%
  \siciteaffiliation{\textsuperscript{1}Universidade Tecnológica Federal do Paraná, Campus Cornélio Procópio}%
}

% O regulamento exige que este e-mail seja o mesmo cadastrado no SISPEQ.
\renewcommand{\sicitedadosemail}{gian.2000@alunos.utfpr.edu.br}

% CONFERIR: o evento tem 19 áreas temáticas. Selecione a correspondente no
% momento da submissão e replique aqui exatamente o número e o nome.
\renewcommand{\sicitedadosarea}{[Conferir número] -- Engenharias}

\renewcommand{\sicitedadospalavraschave}{Comunicação semântica; Aprendizado federado; Redes 6G; Autoencoders; Computação de borda.}
